% Generated draft: requires mathematical and source review.
\documentclass[11pt]{article}
\usepackage[margin=25mm]{geometry}
\usepackage{fontspec}
\setmainfont{Latin Modern Roman}
\usepackage{amsmath,amssymb,mathtools,unicode-math}
\setmathfont{Latin Modern Math}
\usepackage{enumitem,needspace,xurl,hyperref}
\setlength{\parindent}{0pt}
\setlength{\parskip}{5pt}
\setlength{\emergencystretch}{3em}
\newcommand{\N}{\mathbb{N}}
\newcommand{\Z}{\mathbb{Z}}
\newcommand{\Q}{\mathbb{Q}}
\newcommand{\R}{\mathbb{R}}
\newcommand{\C}{\mathbb{C}}
\newcommand{\F}{\mathbb{F}}
\newcommand{\A}{\mathbb{A}}
\newcommand{\PP}{\mathbb{P}}
\newcommand{\E}{\mathbb{E}}
\newcommand{\eps}{\varepsilon}
\newcommand{\dd}{\,\mathrm{d}}
\begin{document}
\section*{UnsolvedMath ID: 700673}
\textbf{Source problem code:} catalogue entry
\subsection{Definitions and mathematical statement}
Under the usual nonempty convention, a formation of finite groups is a nonempty class \(\mathcal F\) of finite groups, closed under isomorphisms, homomorphic images, and finite subdirect products. Here a subdirect product of \(G_1,\ldots,G_r\) is a subgroup of \(G_1\times\cdots\times G_r\) whose projection onto every factor is surjective. The supplied record does not specify the convention concerning the empty class.

Let \((\mathfrak S,\ast)\) denote the intended semigroup of formations of finite groups, and let \(\mathfrak D\subseteq\mathfrak S\) denote the formations described in the source as ``undecomposable.'' The supplied record does not define either the multiplication \(\ast\) or ``undecomposable,'' including any convention concerning identity elements or trivial factors. These are unresolved parts of the formulation; the notation here retains them as supplied rather than selecting definitions.

For any specified interpretation of \(\mathfrak S\), \(\ast\), and \(\mathfrak D\), the subsemigroup generated by \(\mathfrak D\) is
\[
\langle\mathfrak D\rangle
=\{\mathcal F_1\ast\cdots\ast\mathcal F_n:
 n\geq 1,\ \mathcal F_1,\ldots,\mathcal F_n\in\mathfrak D\}.
\]
The question is whether this subsemigroup is free. Interpreting the intended generators as the undecomposable formations themselves, this asks whether, for every pair of positive integers \(n,m\) and all \(\mathcal F_1,\ldots,\mathcal F_n,\mathcal H_1,\ldots,\mathcal H_m\in\mathfrak D\),
\[
\mathcal F_1\ast\cdots\ast\mathcal F_n
=\mathcal H_1\ast\cdots\ast\mathcal H_m
\quad\Longleftrightarrow\quad
n=m\ \text{ and }\ \mathcal F_i=\mathcal H_i
\text{ for every }1\leq i\leq n.
\]
Equivalently, the evaluation map from the free semigroup of finite nonempty words over \(\mathfrak D\) to \(\langle\mathfrak D\rangle\) is injective. This explicit formulation interprets freeness as unique ordered factorization into the stated generators; its applicability requires confirmation of the missing terminology and conventions.
\subsection{Short English statement}
A. N. Skiba asks whether the subsemigroup generated by the undecomposable formations of finite groups is free. With those formations as free generators, this means that two nonempty products are equal exactly when their factors agree in number and order. The supplied record leaves the multiplication and the meaning of ``undecomposable'' undefined, so this remains a qualified statement awaiting mathematical review.
\subsection{Sources}
\begin{itemize}
\item Local UnsolvedMath registry, permanent ID \(700673\), exact source problem code \texttt{catalogue entry}; supplied complete source record, upstream revision \texttt{d371ee917ac9647fbb06cd6fc529deec6f3dda82}.
\item Supplied attribution: A. N. Skiba, proposed in 1986; \emph{The Kourovka Notebook (Unsolved Problems in Group Theory)}, original problem 10.58, 10th issue (1986). The record identifies the 21st issue edition (2026), \texttt{arXiv:1401.0300v48}. The collection description names E. I. Khukhro and V. D. Mazurov as editors.
\item The local record labels the problem ``open.'' This is supplied registry status, not an independently verified assessment. The cited publication has not been consulted for this draft; no external source verification is claimed.
\end{itemize}
\end{document}
